%
% patch_application.tex
%
% A Z specification of the lux display-side patch-application boundary:
% a batch of update patches applied to a scene's element tree by the
% display message loop.  Formalises the crash-freedom discipline hardened
% empirically across FOUR fix/review rounds -- crash on out-of-range
% ValueError, crash on TypeError, non-atomic partial mutation, and crash
% on an unknown-field structural error -- so that ProB can enumerate every
% patch class and either confirm the four invariants or produce the exact
% offending trace.
%
% Source of record:
%   src/punt_lux/scene/patch_applier.py   -- PatchApplier.apply: the batch
%                                            loop; _apply_patch_set: unknown
%                                            field check (raises BEFORE the
%                                            value-rejection try) then
%                                            apply_set under
%                                            try/except (ValueError, TypeError)
%   src/punt_lux/scene/element_walk.py     -- find: locate or None; a missing
%                                            id is warned and skipped
%   src/punt_lux/domain/element_abc.py     -- apply_patch: snapshot + rollback
%                                            (atomic multi-field set)
%   src/punt_lux/scene/manager.py:207      -- apply_update -> applier.apply
%   src/punt_lux/display/server.py:678     -- the display message loop:
%                                            self._scene_manager.apply_update(msg)
%                                            with NO boundary try/except
%
% Type-check:  fuzz -t docs/patch_application.tex
% Model-check (see Section "Verification" for the exact goal predicates):
%   probcli docs/patch_application.tex -model_check -p DEFAULT_SETSIZE 2
%   probcli docs/patch_application.tex -model_check -nodead \
%       -goal "live=terminated" -p DEFAULT_SETSIZE 2
%
% Formatting follows the fuzz house rules: \quad~ for continuation indent
% (never \t1..\t9), schema lines under 80 columns, predicates broken at
% \land/\lor with the operator starting the continuation line.
%

\documentclass[a4paper,10pt,fleqn]{article}
\usepackage[margin=1in]{geometry}
\usepackage{fuzz}
\usepackage[colorlinks=true,linkcolor=blue,citecolor=blue,urlcolor=blue]{hyperref}

\begin{document}

\title{The lux Display Patch-Application Boundary: a Z Specification}
\author{Formal model of the crash-free patch discipline}
\date{July 2026}
\maketitle

% ============================================================
\section{Introduction}
% ============================================================

An agent sends the display a batch of update patches.  Each patch names an
element by id and either removes it or sets some of its fields.  A patch may
be wrong in several ways: the value it sets may be out of range or not a
number, the field it names may not exist on the element, or the id it names
may not exist in the scene.  The display processes the batch in a single
message-loop turn.  The one property that must hold is that \emph{no patch,
however malformed, terminates that loop}.

The same class of defect --- a patch that terminates the display ---
recurred across four fix/review rounds:

\begin{enumerate}
  \item a crash on an out-of-range \texttt{ValueError} from a setter;
  \item a crash on a \texttt{TypeError} from a setter given the wrong type;
  \item a non-atomic partial mutation, where a multi-field patch applied
    some fields before a later field was rejected;
  \item a crash on an unknown-field structural error, where the field-name
    check raises \emph{before} the value-rejection \texttt{try} and the
    exception propagates through \texttt{apply\_update} to the message
    loop and terminates the display.
\end{enumerate}

Rounds~1 and~2 were closed by catching \texttt{(ValueError, TypeError)}
around the setter call; round~3 by an all-or-nothing snapshot/rollback in
\texttt{Element.apply\_patch}; round~4 is \emph{live} --- a test has
confirmed that \texttt{\{"set":\{"bogus":5\}\}} raises out of
\texttt{\_apply\_patch\_set} and kills the loop.  Empirical case-by-case
patching keeps missing cases.  This document models the boundary as a
finite state machine over a bounded batch, so ProB can enumerate every
patch class exhaustively.

% ============================================================
\section{Basic Types}
% ============================================================

\begin{zed}
  [ELEM]
\end{zed}

\texttt{ELEM} is the set of element ids that may appear in the scene or be
named by a patch.  Two ids suffice to exhibit batch-continuation across
distinct targets (one patch rejected, a later patch to another element
applied); the carrier is bounded by ProB's \texttt{DEFAULT\_SETSIZE}.

% ============================================================
\section{Free Types}
% ============================================================

\begin{zed}
  VALID ::= good | bad
\end{zed}

Whether an element's tracked field value satisfies that field's invariant.
\texttt{good}: in range, correct type.  \texttt{bad}: out of range, NaN, or
wrong type --- the value a setter rejects.  The safety goal is that a
\texttt{bad} value is never \emph{installed}.

\begin{zed}
  LIVE ::= alive | terminated
\end{zed}

The display message loop.  \texttt{alive}: still processing messages.
\texttt{terminated}: an unhandled exception escaped the patch path and
killed the loop --- the defect this model exists to forbid.

\begin{zed}
  KIND ::= setGood | setBadVal | setBadType | setUnknown | removeIt
\end{zed}

The five ways a patch acts on its target.  \texttt{setGood}: set a known
field to a valid value.  \texttt{setBadVal}: set a known field to an
out-of-range or NaN value (a setter \texttt{ValueError}).  \texttt{setBadType}:
set a known field to a value of the wrong type (a setter \texttt{TypeError}).
\texttt{setUnknown}: name a field the element does not have (the structural
\texttt{ValueError} raised before the value-rejection \texttt{try}).
\texttt{removeIt}: remove the element.

A patch is a token carrying its kind and its target id:

\begin{zed}
  TOKEN ::= tSetGood \ldata ELEM \rdata | tSetBadVal \ldata ELEM \rdata \\
  \quad~ | tSetBadType \ldata ELEM \rdata | tSetUnknown \ldata ELEM \rdata \\
  \quad~ | tRemove \ldata ELEM \rdata
\end{zed}

Each constructor is an injection \texttt{ELEM $\inj$ TOKEN}, so a token
determines its target and its kind uniquely.  Two total projections
recover them:

\begin{axdef}
  targetOf : TOKEN \fun ELEM \\
  kindOf : TOKEN \fun KIND
\where
  \forall e : ELEM @ \\
  \quad~ targetOf~(tSetGood~e) = e \land targetOf~(tSetBadVal~e) = e \land \\
  \quad~ targetOf~(tSetBadType~e) = e \land targetOf~(tSetUnknown~e) = e \land \\
  \quad~ targetOf~(tRemove~e) = e \\
  \forall e : ELEM @ \\
  \quad~ kindOf~(tSetGood~e) = setGood \land kindOf~(tSetBadVal~e) = setBadVal \land \\
  \quad~ kindOf~(tSetBadType~e) = setBadType \land \\
  \quad~ kindOf~(tSetUnknown~e) = setUnknown \land kindOf~(tRemove~e) = removeIt
\end{axdef}

Because every \texttt{TOKEN} is exactly one construction, these equations
pin \texttt{targetOf} and \texttt{kindOf} to a unique total function apiece.

The batch is a bounded sequence of tokens.  The bound keeps the carrier
finite so ProB can enumerate every ordering:

\begin{axdef}
  maxlen : \nat
\where
  maxlen = 3
\end{axdef}

% ============================================================
\section{State}
% ============================================================

The state is flat: one schema, five components.  Its predicate carries only
\emph{structural} invariants --- that the tracked-value map is defined
exactly on the present elements, that a partially-mutated element is
present, and the length bound on the queue.  These are coherence
constraints, true in every design.  The four safety properties
(Section~\ref{sec:inv}) are deliberately \emph{absent} here, so a violating
state stays reachable and ProB can find it.

\begin{schema}{PatchSys}
  live : LIVE \\
  present : \power ELEM \\
  prim : ELEM \pfun VALID \\
  partial : \power ELEM \\
  queue : \seq TOKEN
\where
  \dom prim = present \\
  partial \subseteq present \\
  \# queue \leq maxlen
\end{schema}

\texttt{present} is the set of element ids currently in the scene.
\texttt{prim} records, for each present element, whether its tracked field
holds a \texttt{good} or \texttt{bad} value.  \texttt{partial} is a
history component: the set of elements a \emph{rejected} patch left partly
mutated.  In a correct design it is always empty; it is the device that
turns the two-state atomicity property into a one-state invariant ProB can
check.  \texttt{queue} is the unprocessed tail of the current batch.

% ============================================================
\section{Initialization}
% ============================================================

\begin{schema}{Init}
  PatchSys'
\where
  live' = alive \\
  present' = ELEM \\
  prim' = (\lambda e : ELEM @ good) \\
  partial' = \emptyset \\
  queue' = \langle \rangle
\end{schema}

A single initial state: the loop is alive, every element is present with a
\texttt{good} value, nothing is partially mutated, and the batch is empty.

% ============================================================
\section{Building a Batch}
% ============================================================

An agent enqueues patches until the batch is full.  Any token may be
enqueued --- including one whose target is not present, which models a
patch for a missing id:

\begin{schema}{Enqueue}
  \Delta PatchSys \\
  p? : TOKEN
\where
  \# queue < maxlen \\
  queue' = queue \cat \langle p? \rangle \\
  live' = live \\
  present' = present \\
  prim' = prim \\
  partial' = partial
\end{schema}

% ============================================================
\section{Applying a Patch}
\label{sec:apply}
% ============================================================

Each step processes the head of the queue.  The five operation schemas
partition every head token by target-presence and kind, so whenever the
loop is alive and the queue is non-empty exactly one is enabled: the
processing of a batch never gets stuck (Section~\ref{sec:inv},
Invariant~BC).

A patch whose target is not present is a no-op --- the tree walk returns no
location, and the applier warns and skips.  This covers both a set and a
remove for a missing id:

\begin{schema}{ApplyMissing}
  \Delta PatchSys
\where
  live = alive \\
  queue \neq \langle \rangle \\
  targetOf~(head~queue) \notin present \\
  queue' = tail~queue \\
  live' = live \\
  present' = present \\
  prim' = prim \\
  partial' = partial
\end{schema}

A remove of a present element detaches it: the id leaves \texttt{present},
its value map entry and any partial mark go with it.

\begin{schema}{ApplyRemove}
  \Delta PatchSys
\where
  live = alive \\
  queue \neq \langle \rangle \\
  targetOf~(head~queue) \in present \\
  kindOf~(head~queue) = removeIt \\
  present' = present \setminus \{ targetOf~(head~queue) \} \\
  prim' = \{ targetOf~(head~queue) \} \ndres prim \\
  partial' = partial \setminus \{ targetOf~(head~queue) \} \\
  queue' = tail~queue \\
  live' = live
\end{schema}

A set of a valid value on a known field installs \texttt{good}:

\begin{schema}{ApplySetGood}
  \Delta PatchSys
\where
  live = alive \\
  queue \neq \langle \rangle \\
  targetOf~(head~queue) \in present \\
  kindOf~(head~queue) = setGood \\
  prim' = prim \oplus \{ targetOf~(head~queue) \mapsto good \} \\
  present' = present \\
  partial' = partial \\
  queue' = tail~queue \\
  live' = live
\end{schema}

A set that a validated setter \emph{rejects} --- an out-of-range or NaN
value (\texttt{ValueError}) or a wrong-typed value (\texttt{TypeError}) ---
is caught around the setter call, logged, and skipped.  Nothing is
installed, so the \texttt{bad} value never reaches \texttt{prim}; nothing is
mutated, so no element is marked \texttt{partial}; the loop stays alive and
the batch continues:

\begin{schema}{ApplyReject}
  \Delta PatchSys
\where
  live = alive \\
  queue \neq \langle \rangle \\
  targetOf~(head~queue) \in present \\
  kindOf~(head~queue) \in \{ setBadVal, setBadType \} \\
  queue' = tail~queue \\
  live' = live \\
  present' = present \\
  prim' = prim \\
  partial' = partial
\end{schema}

A set naming an unknown field is a \emph{structural} error.  The correct
boundary surfaces it --- logs it and/or returns it to the agent --- per
patch, then continues.  It installs nothing, mutates nothing, and, above
all, does not terminate the loop:

\begin{schema}{ApplySurface}
  \Delta PatchSys
\where
  live = alive \\
  queue \neq \langle \rangle \\
  targetOf~(head~queue) \in present \\
  kindOf~(head~queue) = setUnknown \\
  queue' = tail~queue \\
  live' = live \\
  present' = present \\
  prim' = prim \\
  partial' = partial
\end{schema}

In this correct design \texttt{ApplyReject} and \texttt{ApplySurface} have
the same effect --- skip and continue.  They are kept as separate schemas
because they diverge under the buggy variants of
Section~\ref{sec:fidelity}: round~1/2 breaks \texttt{ApplyReject}, round~4
breaks \texttt{ApplySurface}, and the model must be able to break each
independently.

% ============================================================
\section{Invariants}
\label{sec:inv}
% ============================================================

Four properties must hold across all reachable states.  They are stated
here as predicates over \texttt{PatchSys}; none is placed in the
\texttt{PatchSys} schema, because a predicate placed there would be enforced
as a guard on every successor, silently disabling any operation that would
break it and thereby masking the very defect this model exists to catch.

\paragraph{Invariant CF --- crash-freedom.}
The display loop is never terminated:
\[
  live = alive.
\]
This is the load-bearing property.  Every patch class --- a rejected value,
a rejected type, an unknown field, a missing id, a malformed remove --- is
handled at the boundary; none escapes to the loop.

\paragraph{Invariant AT --- atomicity.}
No element is ever left partially mutated by a rejected patch:
\[
  partial = \emptyset.
\]
A rejected multi-field set restores every field it touched, so the ``keeps
its previous value'' claim is literally true.

\paragraph{Invariant VP --- validity-preservation.}
No \texttt{bad} value is ever installed; every present element's tracked
value is \texttt{good}:
\[
  \ran prim \subseteq \{ good \}.
\]

\paragraph{Invariant BC --- batch-continuation.}
A rejected patch never abandons the rest of its batch.  Because the only
thing that halts processing is termination, and termination never occurs
(Invariant~CF), the observable failure is a batch left unfinished with the
loop dead:
\[
  \lnot\, (live = terminated \land queue \neq \langle \rangle).
\]
Each of \texttt{ApplyReject} and \texttt{ApplySurface} reduces the queue by
exactly its head and frames \texttt{live} at \texttt{alive}, so the
successor patch is processed on the next step.

% ============================================================
\section{Verification}
\label{sec:verify}
% ============================================================

Type-checking is a \texttt{fuzz} pass:
\begin{verbatim}
  fuzz -t docs/patch_application.tex
\end{verbatim}

Each safety property is checked by asking ProB whether its negation is
\emph{reachable}.  A goal that is not found means the property holds on
every reachable state:
\begin{verbatim}
  # Invariant CF (must report: goal NOT found)
  probcli docs/patch_application.tex -model_check -nodead \
    -goal "live=terminated" -p DEFAULT_SETSIZE 2

  # Invariant AT (must report: goal NOT found)
  probcli docs/patch_application.tex -model_check -nodead \
    -goal "partial /= {}" -p DEFAULT_SETSIZE 2

  # Invariant VP (must report: goal NOT found)
  probcli docs/patch_application.tex -model_check -nodead \
    -goal "bad : ran(prim)" -p DEFAULT_SETSIZE 2

  # Invariant BC (must report: goal NOT found)
  probcli docs/patch_application.tex -model_check -nodead \
    -goal "live=terminated & queue /= <>" -p DEFAULT_SETSIZE 2
\end{verbatim}

Deadlock-freedom is the state-space search with no goal; the loop is kept
live by \texttt{Enqueue} while the batch has room and by an
\texttt{ApplyX} while it is non-empty, so a correct design never deadlocks:
\begin{verbatim}
  # Deadlock-freedom (must report: no deadlock, no invariant violation)
  probcli docs/patch_application.tex -model_check -p DEFAULT_SETSIZE 2
\end{verbatim}

% ============================================================
\section{Fidelity: reproducing the four defects}
\label{sec:fidelity}
% ============================================================

A model that cannot reproduce the known defects proves nothing.  Each of
the four rounds corresponds to breaking one guard of the correct design;
the correct spec above is what remains once all four are fixed.

\paragraph{Round 1 --- crash on out-of-range \texttt{ValueError}.}
Remove the catch from the value-rejection path: in \texttt{ApplyReject},
for \texttt{setBadVal}, replace the frame \texttt{live' = live} with
\texttt{live' = terminated}.  The reachability search for Invariant~CF then
returns \emph{goal found} --- a batch containing one \texttt{tSetBadVal}
token drives \texttt{live} to \texttt{terminated}.

\paragraph{Round 2 --- crash on \texttt{TypeError}.}
The same edit for \texttt{setBadType}.  A \texttt{tSetBadType} token
reaches \texttt{terminated}: \emph{goal found}.  Rounds~1 and~2 were two
review rounds precisely because value-rejection and type-rejection were
handled by different \texttt{except} clauses; the model shows they are one
case class.

\paragraph{Round 3 --- non-atomic partial mutation.}
Model the missing rollback: in \texttt{ApplyReject}, replace
\texttt{partial' = partial} with
\texttt{partial' = partial $\cup$ \{targetOf($head$~queue)\}} (a
multi-field set applied its first field before the second was rejected, and
nothing restored it).  The loop stays alive, but Invariant~AT's goal
\texttt{partial /= \{\}} is now \emph{found}.

\paragraph{Round 4 --- crash on unknown-field structural error (live).}
This is the current production behaviour.  In \texttt{ApplySurface} replace
\texttt{live' = live} with \texttt{live' = terminated}: the field-name
\texttt{ValueError} is raised before the value-rejection \texttt{try} and
propagates through \texttt{apply\_update} to the message loop.  Invariant~CF's
goal \texttt{live=terminated} is \emph{found}, and so is Invariant~BC's goal
\texttt{live=terminated \& queue /= <>} --- an unknown-field patch not only
kills the loop but strands every patch after it in the batch.

Under the correct design none of these traces exists: \texttt{ApplyReject}
and \texttt{ApplySurface} frame \texttt{live} at \texttt{alive} and
\texttt{partial} unchanged, \texttt{prim} is only ever set to \texttt{good},
and no operation assigns \texttt{terminated}.  The reachability searches for
CF, AT, VP, and BC therefore return \emph{goal not found} for the correct
design and \emph{goal found} --- the traces above --- for the corresponding
buggy variant.  That difference is the evidence that the model detects the
defects the code was changed, round by round, to prevent.

\end{document}
